\documentclass[a5paper, leqno, 10pt]{article}
\usepackage{preamble}

\begin{document}

\newgeometry{
    left=10mm,
    right=10mm,
    top=10mm,
    bottom=10mm,
    includefoot=false
}

\begin{titlepage}
	\thispagestyle{empty}
    \begin{center}
        МГУ имени М.\,В.~Ломоносова

		\bigskip

        Механико-математический факультет

        \vfill

        {\Large \textbf{Отчёт по задаче №1 практикума на ЭВМ}\par \textbf{<<Конечные уравнения>>}}

		\bigskip

		Задача 1{.}10

        \vfill

        \begin{flushright}
            \textit{Никита Пшеничный}\par группа~403
        \end{flushright}

        \vfill

        Москва\\
        2026
    \end{center}
\end{titlepage}

\restoregeometry

\section{Постановка задачи}

Требуется найти все корни уравнения
\[
	\ln(\cos x) + e^{-x} = 0.
\]

\section{Предварительное исследование}

Обозначим $f(x) \coloneqq \ln(\cos x) + e^{-x}$. Функция $f$ определена на множестве
\[
	\bigcup_{k \in \Z}I_k,\quad I_k \coloneqq \left(2\pi k - \frac{\pi}{2}; 2\pi k + \frac{\pi}{2}\right).
\]
Обозначим $a_k \coloneqq 2\pi k - \frac{\pi}{2}$, $b_k \coloneqq 2\pi k + \frac{\pi}{2}$. Так как
\[
	f(2\pi k) = e^{-2\pi k} > 0,\quad \lim_{x \to a_k +}f(x) = \lim_{x \to b_k -}f(x) = -\infty,
\]
то на каждом из интервалов $(a_k; 2\pi k)$ и $(2\pi k; b_k)$, $k \in \Z$, есть корень функции $f$. Докажем, что он единственный. В самом деле,
\begin{gather*}
	f^{\prime}(x) = -\tg x - e^{-x},\\
	\big(e^xf^{\prime}(x)\big)^{\prime} = -e^x\left(\tg x + \frac{1}{\cos^2x}\right) = -e^x\left(\left(\tg x + \frac{1}{2}\right)^2 + \frac{3}{4}\right) < 0.
\end{gather*}

Следовательно, на каждом из интервалов $I_k$, $k \in \Z$, функция $e^xf^{\prime}(x)$ убывает, причём
\[
	\lim_{x \to a_k +}e^xf^{\prime}(x) = +\infty,\quad \lim_{x \to a_k -}e^xf^{\prime}(x) = -\infty.
\]
Значит, $f^{\prime}$ имеет на каждом интервале $I_k$ ровно один нуль.

Таким образом, на каждом из интервалов $I_k$, $k \in \Z$, функция $f$ выпукла вниз, а значит, имеет не более двух корней на нём. Обозначим корни функции $f$, принадлежащие интервалам $(a_k; 2\pi k)$ и $(2\pi k; b_k)$, через $x_k^{-}$ и $x_k^{+}$ соответственно.

\subsection{Корни уравнения $f(x) = 0$ на $I_k$ при $k \leqslant -1$}

Так как $x_k^{-} - a_k, b_k - x_k^{+} \in \big(0; \frac{\pi}{2}\big)$ и $f(x_k^{-}) = f(x_k^{+}) = 0$, то
\begin{gather*}
	\frac{2(x_k^{-} - a_k)}{\pi} \leqslant \sin(x_k^{-} - a_k) = \cos x_k^{-} \leqslant \exp(-e^{-b_k}),\\
	\frac{2(b_k - x_k^{+})}{\pi} \leqslant \sin(b_k - x_k^{+}) = \cos x_k^{+} \leqslant \exp(-e^{-b_k}).
\end{gather*}
Отсюда $0 < x_k^{-} - a_k,\,b_k - x_k^{+} \leqslant \frac{\pi}{2}\exp(-e^{-b_k})$. При $k \leqslant -1$
\[
	x_k^{-},\,x_k^{+} < b_k = 2\pi k + \frac{\pi}{2} \leqslant -\frac{3\pi}{2},
\]
следовательно, $\lvert x_k^{-}\rvert,\,\lvert x_k^{+}\rvert \geqslant \frac{3\pi}{2}$, а значит,
\[
	\frac{\lvert x_k^{-} - a_k\rvert}{\lvert x_k^{-}\rvert},\,
	\frac{\lvert x_k^{+} - b_k\rvert}{\lvert x_k^{+}\rvert} \leqslant \frac{\frac{\pi}{2}\exp(-e^{3\pi / 2})}{\frac{3\pi}{2}} = \frac{1}{3}\exp(-e^{3\pi / 2}) \approx 1{.}51 \cdot 10^{-49}.
\]

Так как полученная относительная погрешность меньше машинной точности, то при численном решении можно считать, что при $k \leqslant -1$
\[
	x_k^{-} \approx a_k = 2\pi k - \frac{\pi}{2},\quad
	x_k^{+} \approx b_k = 2\pi k + \frac{\pi}{2}.
\]

\subsection{Корни уравнения $f(x) = 0$ на $I_k$ при $k \to \infty$}

Обозначим $L(t) \coloneqq -\ln\cos t$.

\begin{lemma}
	Если $0 < t < \frac{\pi}{2}$, то
	\begin{equation} \label{eq:L_estimates}
		\frac{t^2}{2} \leqslant L(t) \leqslant \frac{t^2}{2\cos^2t}.
	\end{equation}
\end{lemma}

\begin{proof}
	Заметим, что
	\[
		-\ln\cos t = \int_0^t\tg u\diff u.
	\]
	Но $\tg u \geqslant u$ при $0 < u < \frac{\pi}{2}$, следовательно,
	\[
		L(t) \geqslant \int_0^tu\diff u = \frac{t^2}{2}.
	\]
	С другой стороны,
	\[
		\tg u = \int_0^u\frac{1}{\cos^2v}\diff v.
	\]
	Если $0 < v \leqslant t < \frac{\pi}{2}$, то $\frac{1}{\cos^2v} < \frac{1}{\cos^2t}$, следовательно, при $0 < u < t$
	\[
		\tg u \leqslant \frac{1}{\cos^2t}\int_0^u\diff u = \frac{u}{\cos^2t}.
	\]
	Отсюда
	\[
		L(t) \leqslant \frac{1}{\cos^2 t}\int_0^tu\diff u = \frac{t^2}{2\cos^2t},
	\]
	что и требовалось.
\end{proof}

Обозначим $\alpha_k \coloneqq 2\pi k - x_k^{-}$ и $h_k \coloneq \sqrt{2}e^{-\pi k}$, $k \in \Z$, тогда
\[
	L(\alpha_k) = -\ln\cos x_k^{-} = e^{\alpha_k - 2\pi k} = \frac{h_k^2}{2}e^{\alpha_k}.
\]
Из $0 < \alpha_k < \frac{\pi}{2}$ и первого неравенства в \eqref{eq:L_estimates} вытекает, что
\[
	0 < \alpha_k \leqslant h_ke^{\alpha_k / 2} = \sqrt{2}e^{\alpha_k / 2 -\pi k} \leqslant \sqrt{2}e^{-3\pi / 4} \approx 0{.}134 < 1.
\]
Как известно, $e^{t / 2} \leqslant 1 + t$ при $0 \leqslant t \leqslant 1$. Отсюда и из $\alpha_k \leqslant h_ke^{\alpha_k / 2}$ получаем
\begin{equation} \label{eq:upper_estimate}
	\alpha_k \leqslant \frac{h_k}{1 - h_k}.
\end{equation}
Из второго неравенства в \eqref{eq:L_estimates} следует, что
\[
	\alpha_k \geqslant h_k\cos\alpha_ke^{\alpha_k / 2}.
\]
Так как при $0 \leqslant t \leqslant 1$
\[
	e^{t / 2} \geqslant 1,\quad \cos t \geqslant 1 - \frac{t^2}{2} \geqslant 1 - t,
\]
то для $0 < \alpha_k < 1$ получаем
\begin{equation} \label{eq:lower_estimate}
	\alpha_k \geqslant \frac{h_k}{1 + h_k}.
\end{equation}
Объединяя оценки \eqref{eq:upper_estimate} и \eqref{eq:lower_estimate}, получаем
\[
	\frac{h_k}{1 + h_k} \leqslant \alpha_k \leqslant \frac{h_k}{1 - h_k}.
\]
Отметим также, что при $k \geqslant 1$
\[
	h_k \leqslant \sqrt{2}e^{-\pi} \approx 0{.}061 < \frac{1}{3}.
\]
Следовательно,
\[
	\lvert\alpha_k - h_k\rvert \leqslant \frac{h_k}{1 - h_k} - h_k = \frac{h_k^2}{1 - h_k} < \frac{3h_k^2}{2} = 3e^{-2\pi k}.
\]
Аналогично для $\beta_k \coloneqq x_k^{+} - 2\pi k$ можно получить оценку
\[
	\lvert\beta_k - h_k\rvert \leqslant 3e^{-2\pi k}.
\]
Таким образом,
\[
	\lvert x_k^{-} - (2\pi k - h_k)\rvert,\,\lvert x_k^{+} - (2\pi k + h_k)\rvert \leqslant 3e^{-2\pi k}.
\]
Относительные погрешности равны
\begin{equation} \label{eq:final_relative_estimate}
	\frac{|x_k^{-} - (2\pi k - h_k)|}{|x_k^{-}|},\,\frac{|x_k^{+} - (2\pi k + h_k)|}{|x_k^{+}|} < \frac{3e^{-2\pi k}}{2\pi k - \frac{\pi}{2}} \to 0,\ k \to \infty.
\end{equation}

\subsection{Численное решение уравнения $f(x) = 0$}

Пусть $\eps > 0$ --- заданная относительная погрешность решения. Как было показано выше, при $k \leqslant -1$
\[
	\frac{\big|x_k^{-} - \big(2\pi k - \frac{\pi}{2}\big)\big|}{|x_k^{-}|},\,
	\frac{\big|x_k^{+} - \big(2\pi k + \frac{\pi}{2}\big)\big|}{|x_k^{+}|} < \eps
\]
для всех допустимых $\eps$ (минимальная допустимая). Поэтому для $k \leqslant -1$ в качестве корней уравнения $f(x) = 0$ на $I_k$ будут использоваться значения $2\pi k \pm \frac{\pi}{2}$. Отметим, что в этих точках функция $f$ не определена, поэтому проверять корни подстановкой некорректно.

На каждом интервале $I_k$, $k \geqslant 0$, будем вычислять корни $x_k^{-}, x_k^{+} \in I_k$ уравнения $f(x) = 0$ методом деления отрезка пополам. При $k \geqslant 1$ выполнены оценки \eqref{eq:final_relative_estimate}, поэтому перед вычислением корней на интервале $I_k$ будем проверять, что
\[
	\frac{3e^{-2\pi k}}{2\pi k - \frac{\pi}{2}} \geqslant \eps.
\]
В противном случае в качестве корней будем использовать величины $2\pi k \pm h_k$. Отметим, что уже при $k = 5$ относительная погрешность корней не превышает
\[
	\frac{3e^{-10\pi}}{10\pi - \frac{\pi}{2}} \approx 2{.}28 \cdot 10^{-15} < 10^{-12}.
\]

\section{Результаты работы программы}

Приведём результаты работы программы для нескольких значений относительной погрешности $\eps$. При $\eps = 10^{-2}$ программа вычисляет корни только при $k = 0$:

\begin{center}
	\begin{tabular}{c|cc}
		$k$ & $x_k^-$                 & $x_k^+$                 \\[.5mm]
		\hline
		0   & $-1{.}5585244804918115$ & $0{.}85289331806441648$ 
	\end{tabular}
\end{center}

\noindent
Начиная с $k = 1$, используется асимптотическая формула
\[
	x_k^{\pm} \approx 2\pi k \pm \sqrt{2}e^{-\pi k},
\]
для которой относительная погрешность не превосходит $1{.}19 \cdot 10^{-3} < 10^{-2}$.

При $\eps = 10^{-5}$ программа вычисляет корни при $k = 0, 1$:

\begin{center}
	\begin{tabular}{c|cc}
		$k$ & $x_k^-$               & $x_k^+$               \\[.5mm]
		\hline
		0   & $-1.5623234797868095$ & $0.85754918944014769$ \\
		1   & $ 6.2201482841774105$ & $6.3424832520112906$
	\end{tabular}
\end{center}

Начиная с $k = 2$, используется асимптотическая формула; её относительная погрешность не превосходит $9{.}52 \cdot 10^{-7} < 10^{-5}$.

При $\eps = 10^{-8}$ программа вычисляет корни при $k = 0, 1, 2$:

\begin{center}
	\begin{tabular}{c|cc}
		$k$ & $x_k^-$               & $x_k^+$               \\[.5mm]
		\hline
		0   & $-1.5623145501348534$ & $0.85754930062192170$ \\
		1   & $ 6.2201352232447986$ & $6.3424958448101165$  \\
		2   & $12.563726126605642$  & $12.569008173732678$
	\end{tabular}
\end{center}

Начиная с $k = 3$, используется асимптотическая формула; её относительная погрешность не превосходит $1{.}13 \cdot 10^{-9} < 10^{-8}$.

При $\eps = 10^{-12}$ программа вычисляет корни при $k = 0, 1, 2, 3, 4$:

\begin{center}
	\begin{tabular}{c|cc}
		$k$ & $x_k^-$               & $x_k^+$               \\[.5mm]
		\hline
		0   & $-1.5623145607652962$ & $0.85754930157839016$ \\
		1   & $ 6.2201352007238606$ & $6.3424958806230372$  \\
		2   & $12.563726158801270$  & $12.569008095226552$  \\
		3   & $18.849441788681808$  & $18.849670041389452$  \\
		4   & $25.132736296878630$  & $25.132746160558060$
	\end{tabular}
\end{center}

Начиная с $k = 5$, используется асимптотическая формула; её относительная погрешность не превосходит $2{.}28 \cdot 10^{-15} < 10^{-12}$.

\end{document}

